
% ======================================================================
\chapter{Introduzione}
\label{sec:intro}
% ======================================================================

Un linguaggio di programmazione è \textbf{reversibile} quando ogni sua
computazione può essere ripercorsa all'indietro fino a ritrovare esattamente lo
stato di partenza. Non si tratta di registrare una traccia da cui ricostruire il
passato: in un linguaggio reversibile l'informazione necessaria a tornare
indietro è già contenuta nello stato corrente, perché ogni costrutto è
progettato per non distruggerla. L'assegnamento distruttivo $x := e$, che
cancella il valore precedente di $x$, non appartiene a un linguaggio di questo
tipo; al suo posto compaiono aggiornamenti invertibili come $x \mathrel{+}= e$,
il cui inverso $x \mathrel{-}= e$ è immediato.

L'interesse per questo modo di calcolare ha tre radici indipendenti. La prima è
fisica: il principio di Landauer lega la cancellazione di un bit a una
dissipazione minima di energia, e un calcolo che non cancella può in linea di
principio evitarla. La seconda è pragmatica: molti problemi si presentano
naturalmente in coppie di trasformazioni inverse (compressione e
decompressione, codifica e decodifica, serializzazione e parsing) e un
linguaggio reversibile permette di scrivere \emph{un solo} programma invece di
due, eliminando per costruzione la possibilità che le due metà divergano. La
terza riguarda l'affidabilità: poter tornare indietro è il meccanismo di base
del \emph{recovery}, e infatti il \emph{rollback} di una transazione non è altro
che esecuzione all'indietro fino a un punto sicuro.

\medskip
\noindent
È proprio la terza radice a rendere interessante il caso \textbf{concorrente},
ed è qui che si colloca questo lavoro. In un sistema sequenziale l'esecuzione
all'indietro è una nozione semplice, perché c'è un solo passato da ripercorrere.
In presenza di più thread il passato non è più lineare: due rami che procedono
in parallelo possono aver interagito, e disfare un passo di uno può richiedere
di disfare prima i passi dell'altro che ne dipendono. Il capostipite dei
linguaggi reversibili, \textbf{Janus}~\cite{YokoyamaGluck2007}, è puramente
sequenziale; la reversibilità in contesto concorrente è stata studiata
soprattutto su calcoli di processi~\cite{LaneseMezzinaSchmittStefani2011,%
LaneseLienhardtMezzinaSchmittStefani2013}, che sono modelli e non linguaggi di
programmazione. Il territorio in mezzo (un linguaggio imperativo strutturato,
con la sintassi e i costrutti che un programmatore si aspetta, che sia insieme
reversibile e concorrente) è poco popolato: i tentativi esistono, ma finora
hanno definito i costrutti concorrenti attraverso la loro traduzione, senza una
semantica formale su cui poggiare le proprietà che ci si aspetta da un
linguaggio reversibile.

\medskip
\noindent
\textbf{Kairos} sta in quel territorio. Il nome è il greco
\kairosgr{} (\emph{kairós}): non il tempo che scorre, che è \chronosgr{}
(\emph{chrónos}), ma il momento opportuno, l'istante in cui una cosa va fatta. La distinzione è
pertinente a un linguaggio reversibile, dove il tempo non è una freccia ma una
coordinata percorribile nei due versi, e ciò che conta di un passo non è
\emph{quando} avviene ma che sia quello giusto: l'inversa di un programma non ne
riavvolge la storia, ne esegue un altro che riporta allo stato di partenza.

Al nucleo di Janus (assegnamenti invertibili, condizionali a doppia guardia,
cicli \texttt{from\,\dots\,until}, procedure invocabili con \texttt{call} e
\texttt{uncall}, interi e array) Kairos aggiunge il parallelismo esplicito
\texttt{par\,\dots\,and\,\dots\,rap}, la comunicazione sincrona su canali
\texttt{ssend}/\texttt{srecv} con protocolli ricostruiti dal programma, uno stack come struttura dati primitiva
e le transazioni \texttt{try\,\dots\,rollback\,\dots\,yrt}.
La scelta di progetto che tiene insieme il tutto è la \emph{disgiunzione}: rami
paralleli operano su porzioni disgiunte dello store e comunicano solo per
canali. Da questa singola ipotesi discende la confluenza degli interleaving che
non toccano canali (\S\ref{sec:par-conf}), e quindi il fatto che il
non-determinismo residuo non sia osservabile sullo store finale: la
concorrenza di Kairos serve a esprimere il parallelismo dei problemi, non a
introdurre corse critiche.

\medskip
\noindent
Il contributo di questo documento è la semantica formale del linguaggio.
Diamo una semantica operazionale strutturale in stile small-step, definita come
\textbf{estensione} della semantica di Janus di Lami, Lanese e
Stefani~\cite{LamiLaneseStefani2024}: ereditiamo da quella base la forma dei
giudizi, i contesti di valutazione, la configurazione terminale e il
trattamento degli errori, e aggiungiamo le regole per i costrutti che Kairos
introduce. Definiamo poi l'operatore di inversione $\Inv{\cdot}$, sintattico e
calcolabile in tempo lineare, e mostriamo in che senso realizza l'esecuzione
all'indietro.

\section{Notazione e convenzioni}
\label{sec:notazione}

I giudizi hanno la forma $\cfg{c}{\sigma} \to \cfg{c'}{\sigma'}$, con lo store a
sinistra come nella semantica di base, non il turnstile big-step
$\sigma \vdash e \Downarrow v$ usuale nella letteratura su Janus. Non è una
variante stilistica: interleaving dei rami e sincronizzazione sui canali
(\cref{sec:par}) si esprimono molto più naturalmente passo per passo, perché un
big-step collassa l'intera esecuzione di un comando in un unico giudizio e i
punti in cui intercalare l'altro ramo restano impliciti. È la stessa motivazione
dichiarata in~\cite{LamiLaneseStefani2024}, che adotta lo small-step per Janus
proprio in vista di una futura estensione concorrente.

Una convenzione tipografica riguarda ciò che sta \emph{fuori} dalla relazione di
transizione. Le regole d'errore, che producono $\err$, portano il nome in
\textcolor{red!60!black}{rosso}. I vincoli di buona formazione sono distinti
secondo il momento in cui li si verifica: \static{in verde} quelli controllati
sul testo del programma, prima di eseguirlo; \dynamic{in arancio} quelli
sorvegliati durante l'esecuzione, perché riguardano stato che il testo del
programma non determina.

\section{Struttura della tesi}

Il \cref{chap:background} richiama Janus, la semantica small-step su cui ci
basiamo e i lavori sulla reversibilità concorrente. Il \cref{chap:semantica} è
il nucleo tecnico: sintassi astratta, sistema di transizione, il confine fra ciò
che è ereditato dalla base e ciò che Kairos aggiunge, e le regole di riduzione
dagli assegnamenti fino a \texttt{par} e alle transazioni, con gli errori a
runtime. Il \cref{chap:reversibilita} definisce l'operatore di inversione, ne
dimostra le proprietà e riassume in un quadro unico tutte le regole. Seguono il
\cref{chap:programmare} sull'uso del linguaggio, il \cref{chap:interprete}
sull'implementazione, il \cref{chap:correlati} sul confronto con i lavori
esistenti e il \cref{chap:conclusioni} con le conclusioni.

% ======================================================================
\chapter{Background}
\label{chap:background}
\label{sec:back}
% ======================================================================

\subsubsection{Janus e la programmazione reversibile}

\textbf{Janus} è stato definito intorno al 1982 da Lutz e Derby al
Caltech~\cite{LutzDerby1982}, per curiosità, e riportato all'attenzione della comunità da Yokoyama e
Glück~\cite{YokoyamaGluck2007}, che ne hanno dato la prima semantica formale e
un self-interpreter invertibile. È un linguaggio imperativo strutturato in cui
ogni costrutto è invertibile per costruzione, secondo due idee ricorrenti.

La prima è l'\textbf{aggiornamento invertibile}. Al posto dell'assegnamento
distruttivo Janus ha solo aggiornamenti della forma $x \mathrel{\odot}= e$ con
$\odot \in \{+,-,\hat{\ }\}$, dove $x$ non compare in $e$: il vecchio valore di
$x$ è recuperabile dal nuovo applicando l'operazione inversa, e la condizione
sull'occorrenza di $x$ garantisce che il valore di $e$ sia lo stesso prima e
dopo il passo.

La seconda è la \textbf{doppia guardia}. Un condizionale irreversibile
$\mathtt{if}\ e\ \mathtt{then}\ c_1\ \mathtt{else}\ c_2$ non è invertibile
perché, arrivati alla fine, non si sa più quale ramo sia stato preso. Janus
richiede allora una seconda espressione, l'\emph{asserzione} in coda:
\[
\mathtt{if}\ e_1\ \mathtt{then}\ c_1\ \mathtt{else}\ c_2\ \mathtt{fi}\ e_2 ,
\]
con l'obbligo che $e_2$ valga vero all'uscita se e solo se si è preso il ramo
\texttt{then}. All'indietro è $e_2$ a decidere il ramo ed $e_1$ a fare da
asserzione: il ruolo delle due espressioni si scambia. Lo stesso schema governa
il ciclo
\[
\mathtt{from}\ e_1\ \mathtt{do}\ c_1\ \mathtt{loop}\ c_2\ \mathtt{until}\ e_2 ,
\]
dove $e_1$ deve valere solo alla prima iterazione ed $e_2$ solo all'ultima, così
che la condizione di ingresso all'indietro sia $e_2$ e quella di uscita $e_1$.
I due corpi sono eseguiti secondo la traccia $c_1\,[\,c_2\ c_1\,]^{*}$: $c_1$
porta avanti il lavoro dell'iterazione, $c_2$ fa avanzare il contatore fra
un'iterazione e la successiva. Kairos adotta la stessa forma, con le clausole
$\mathtt{do}$ e $\mathtt{loop}$ facoltative: un corpo assente vale
$\mathtt{skip}$.

Su queste basi Janus è \emph{r-Turing completo}: calcola esattamente le funzioni
iniettive computabili, senza generare spazzatura. Le sue procedure si invocano
con \texttt{call} in avanti e \texttt{uncall} all'indietro, e l'inverso di un
programma si ottiene per trasformazione puramente sintattica.

\subsubsection{Una semantica small-step per Janus}

Le presentazioni classiche di Janus, a partire da~\cite{YokoyamaGluck2007},
usano una semantica \emph{big-step}: un giudizio $\sigma \vdash c \Downarrow
\sigma'$ mette in relazione lo store prima e dopo l'esecuzione \emph{completa}
di $c$. È una scelta comoda per un linguaggio sequenziale, ma collassa in un
unico passo l'intera computazione, e con essa i punti in cui un'altra
esecuzione potrebbe intercalarsi.

Lami, Lanese e Stefani~\cite{LamiLaneseStefani2024} hanno dato una semantica
\textbf{small-step} di Janus, con giudizi della forma $\langle \sigma, c\rangle
\to \langle \sigma', c'\rangle$ fra configurazioni, contesti di valutazione per
espressioni e comandi, configurazione terminale $\langle \sigma,
\mathtt{skip}\rangle$ ed errori segnalati da un esito distinto $\bot$. Uno degli
scopi dichiarati di quel lavoro è proprio predisporre il terreno per estensioni
concorrenti, sulle quali le proprietà di interesse si definiscono più
agevolmente in forma small-step. Kairos parte da lì: adottarne le convenzioni
significa che le scelte di presentazione sono già state fatte e argomentate a
monte, e che le regole nuove sono solo quelle dei costrutti nuovi. La
\S\ref{sec:base} elenca esattamente il confine fra le due cose.

Una semantica alternativa per Janus, in stile \emph{reversible}, è quella di
Lanese e Vidal~\cite{LaneseVidal2026}, in cui le regole in avanti e quelle
all'indietro sono date insieme sullo stesso sistema di transizione anziché
derivare le seconde da un inverter sintattico.

\subsubsection{Reversibilità e concorrenza}

Fuori dai linguaggi imperativi, la reversibilità in presenza di concorrenza è
stata studiata a fondo sui \emph{calcoli di processi}. Il punto delicato è che
in un sistema concorrente non esiste un unico ordine di esecuzione da
rovesciare: la nozione corretta è la \emph{reversibilità causalmente
consistente}, per cui un passo può essere disfatto solo dopo che sono stati
disfatti tutti i passi che da esso dipendono causalmente. In quest'ottica la
reversibilità va anche \emph{controllata}, perché tornare indietro senza limiti
non è utile a nulla: in $\rho\pi$ e in
croll-$\pi$~\cite{LaneseMezzinaSchmittStefani2011,%
LaneseLienhardtMezzinaSchmittStefani2013} il ritorno all'indietro è governato da
primitive esplicite, e il meccanismo che ne risulta (si torna indietro fino a
un punto scelto e da lì si prosegue diversamente) è esattamente quello delle
transazioni con compensazione. Le transazioni
\texttt{try\,\dots\,rollback\,\dots\,yrt} di Kairos (\S\ref{sec:try}) sono la
versione a linguaggio di programmazione della stessa idea.

Kairos adotta però una strada più restrittiva, e per questo più semplice: rami
paralleli lavorano su porzioni \emph{disgiunte} dello store e interagiscono
soltanto per comunicazione sincrona. Non c'è memoria condivisa, quindi non c'è
dipendenza causale da tracciare fra passi che non comunicano, e la consistenza
causale si riduce alla confluenza dimostrata in \S\ref{sec:par-conf}. Il prezzo
è che alcuni schemi di programmazione a memoria condivisa non sono esprimibili;
il guadagno è che un programma che rispetta la disgiunzione non ha corse
critiche.

Fra i calcoli di processo e i linguaggi di programmazione esiste comunque un
precedente diretto: Janus è già stato esteso con un costrutto di composizione
parallela e uno scambio di messaggi, e l'estensione è stata realizzata con un
compilatore verso C++~\cite{Sammarchi2019}. Quel lavoro definisce i costrutti
attraverso il codice in cui vengono tradotti, non attraverso regole di
riduzione, e mantiene uno stato condiviso fra i rami; è quindi il punto di
partenza più vicino a quello di questa tesi, e il \cref{chap:correlati} lo
riprende per un confronto puntuale.

\subsubsection{Canali e tipi di sessione}

La comunicazione di Kairos avviene su canali \emph{sincroni}, con protocolli
ricostruiti dal programma,
con il vincolo che ogni canale abbia in ogni istante al più due titolari, una
\emph{sessione binaria}. È una restrizione che viene dalla teoria dei
\textbf{tipi di sessione}~\cite{HondaVasconcelosKubo1998,%
HondaYoshidaCarbone2008}, dove il tipo di un canale non descrive solo i dati
scambiati ma la \emph{sequenza} degli scambi, e i tipi dei due estremi devono
essere duali. In Kairos il vincolo di sessione binaria serve a garantire che il
partner di ogni comunicazione sia determinato: senza
di esso un $\mathtt{ssend}$ potrebbe sincronizzarsi con più $\mathtt{srecv}$
diversi, e la scelta sarebbe un non-determinismo osservabile sullo store, in
contrasto con la confluenza.

Il vincolo ha due forme, e la scelta fra le due non è di dettaglio. Nella forma
\emph{statica} un canale appartiene agli stessi due rami per tutta la sua vita,
e questo si legge dal testo del programma. Nella forma \emph{dinamica} un canale
può stare prima fra due rami e poi fra altri due, purché il primo titolare
smetta di usarlo prima che il nuovo cominci: è ciò che nella teoria dei tipi di
sessione si chiama \textbf{delega}, il meccanismo con cui, in stile
$\pi$-calcolo, un processo passa a un terzo il canale su cui sta parlando. La
seconda è strettamente più espressiva, e Kairos adotta quella: il vincolo è
sorvegliato osservando l'esecuzione (\S\ref{sec:chan}), non leggendo il testo del
programma, perché il testo non basta a seguire un canale che cambia titolare.

\medskip
\noindent
Va detto dove sta il confine di questo lavoro. Nella semantica il vincolo di
sessione binaria e l'ipotesi di disgiunzione sono \emph{assunzioni} sotto cui
vale la confluenza, non un sistema di tipi: le regole delle
\S\S\ref{sec:esp}--\ref{sec:cmd} non impediscono di scrivere un programma con
una corsa critica, si limitano a non garantire nulla su di esso. Il compilatore
ne applica una verifica conservativa (nessun intero o array scritto in un ramo
e usato in un altro, anche tramite chiamate; nessuno stack condiviso), e
l'interprete segnala, quando se ne accorge, due aggiornamenti simultanei della
stessa variabile intera, senza che questo sia una garanzia; un sistema di tipi
completo resta da fare. La separazione è
deliberata e segue la pratica corrente, in cui il calcolo si definisce senza
preoccuparsi delle corse e un sistema di tipi sovrapposto le esclude; i tipi di
sessione, binari o multiparty, sono precisamente una famiglia di sistemi di tipi
di questo genere. In \S\ref{sec:chan} diamo la disciplina binaria per i canali di
Kairos, delega compresa, e i tipi di sessione che ne descrivono il protocollo:
i tipi non si dichiarano ma si ricostruiscono dal programma, e il compilatore
rifiuta due estremi di un canale con protocolli in conflitto dimostrabile. Restano fuori la scelta etichettata e i canali
con più di due partecipanti.

% ======================================================================
\chapter{Kairos: sintassi e semantica operazionale}
\label{chap:semantica}

\section{Sintassi astratta}
\label{sec:sintassi}
% ======================================================================

Insiemi di base: numeri interi $m \in \mathbb{Z}$, variabili $v \in \mathit{Var}$.
Categorie derivate:

\paragraph{Espressioni aritmetiche $e \in \mathit{Exp}$.}
\begin{align*}
e ::= m \mid v \mid a[e] \mid e \oplus e
\qquad
\oplus \in \{+,\,-,\,*,\,/,\,\%\}
\end{align*}

\noindent
$a[e]$ è la lettura della cella di indice $e$ dell'array $a$. Gli operatori
moltiplicativi compaiono solo dentro le espressioni, mai come assegnamento:
$*$ e $/$ non hanno un'inversa che restituisca l'operando, mentre a destra di
$\mathbin{+\!\!=}$ non tolgono nulla, perché l'inversa dell'assegnamento è
$\mathbin{-\!\!=}$ con la \emph{stessa} espressione, e questa dipende da
variabili che il comando non tocca. $/$ e $\%$ sono la divisione intera e il
resto; il divisore nullo non ha risultato e produce un errore:
\[
\dfrac{}{\cfg{m / 0}{\sigma} \toe \err}\ \rnameerr{Div-Err}
\]
con la regola gemella per $\%$. È un fallimento come gli altri
(\S\ref{sec:err}), non un valore convenzionale: un valore convenzionale
renderebbe l'operazione non iniettiva proprio dove serve che lo sia.
Nel testo dei programmi i letterali sono non negativi (un valore negativo si
scrive $0 - m$), salvo i valori di $\mathtt{local}$ e $\mathtt{delocal}$;
l'interprete rappresenta gli interi con segno a 64 bit in complemento a due.
Le operazioni sono quindi modulo $2^{64}$: un risultato che trabocca si
riavvolge (anche $-2^{63} / -1$, che vale $-2^{63}$), e un letterale si legge
modulo $2^{64}$, così che si scrivono anche le costanti senza segno fino a
$2^{64}-1$; un letterale più grande è un errore di compilazione.

\paragraph{Condizioni $b \in \mathit{Cond}$.}
\begin{align*}
b ::= e = e \mid e \neq e \mid e \geq e \mid e \leq e \mid e > e \mid e < e
  \mid b \land b \mid b \lor b
\end{align*}

\noindent
I connettivi $\land$ e $\lor$ sono quelli di Janus: una guardia può combinare
più confronti, come in $(t = 0) \land (\mathit{num} > 1)$. Nell'implementazione confronti e connettivi sono a
loro volta espressioni che valgono $1$ o $0$, e una guardia è vera quando è
diversa da zero: si può quindi scrivere anche $x \mathbin{+\!\!=} (a < b)$. I
connettivi non cortocircuitano: entrambi gli operandi sono sempre valutati, e
ciascuno può produrre \textsc{Idx-Err} o \textsc{Div-Err}.

\paragraph{Comandi $c \in \mathit{Com}$.}
\begin{align*}
c ::={}& \mathtt{skip} && \text{(comando vuoto)}\\
   \mid{}& \ell \mathbin{+\!\!=} e \mid \ell \mathbin{-\!\!=} e \mid \ell \mathbin{\hat{}\!=} e \mid \ell \mathbin{<\!=\!>} \ell && \text{(assegnamenti reversibili)}\\
   \mid{}& \mathtt{local}\ v = m \mid \mathtt{delocal}\ v = m && \text{(scope)}\\
   \mid{}& \mathtt{local}\ a[n] = m \mid \mathtt{delocal}\ a[n] = m && \text{(scope, array)}\\
   \mid{}& c\ ;\ c && \text{(sequenza)}\\
   \mid{}& \mathtt{if}\ b\ [\mathtt{then}\ c]\ [\mathtt{else}\ c]\ \mathtt{fi}\ b' && \text{(condizionale)}\\
   \mid{}& \begin{array}[t]{@{}l@{}}\mathtt{from}\ b_1\ [\mathtt{do}\ c_1]\\ \quad[\mathtt{loop}\ c_2]\ \mathtt{until}\ b_2\end{array} && \text{(ciclo)}\\
   \mid{}& \mathtt{call}\ f(\bar v) \mid \mathtt{uncall}\ f(\bar v) && \text{(chiamata / inversa)}\\
   \mid{}& \mathtt{push}(v,s) \mid \mathtt{pop}(v,s) && \text{(stack)}\\
   \mid{}& \mathtt{ssend}(\langle \bar v \rangle, \mathit{ch}) \mid \mathtt{srecv}(\langle \bar v \rangle, \mathit{